\documentclass[10pt,a4paper]{amsart}
\usepackage[margin=19mm]{geometry}
\usepackage{amsmath,amssymb,mathtools}
\usepackage[scr=rsfs,cal=euler]{mathalfa}
\usepackage{xfrac}
\usepackage[unicode,hidelinks,bookmarks=false]{hyperref}
\newcommand{\ie}{i.e.\ }
\newcommand{\cf}{cf.\ }
\newcommand{\resp}{respectively\ }
\newcommand{\Subsection}{\subsection}
\providecommand{\nobreakdash}{\nobreak\hbox{-}}
\setlength{\emergencystretch}{2em}
\multlinegap0pt
\usepackage{tikz}
\usetikzlibrary{cd}
\usepackage{mathtools}
\usepackage{dsfont}
\usepackage{amssymb}
\usepackage{tensor}
\usepackage{tikz}
\usepackage{xcolor}
\newtheorem{lemma}{Lemma}
\DeclareMathOperator{\For}{For}
\newcommand{\Hilb}{\mathrm{Hilb}}
\DeclareMathOperator{\Vect}{Vec}
\newcommand{\bB}{\mathbb B}
\newcommand{\cC}{\mathcal C}
\newcommand{\cL}{\mathcal L}
\newcommand{\id}{\mathrm{id}}
\title{Monadic reconstruction of unitary Drinfeld centers and Factorization Homology}
\author{Lucas Hataishi}
\date{Source: arXiv:2512.08848v2 (CC BY 4.0)}
\begin{document}
\maketitle

\noindent\emph{Source and licence.} Monadic reconstruction of unitary Drinfeld centers and Factorization Homology. Version arXiv:2512.08848v2 (CC BY 4.0), distributed by arXiv under the Creative Commons Attribution 4.0 International licence, http://creativecommons.org/licenses/by/4.0/, as recorded in the arXiv licence metadata for this exact version and shown on the abstract page https://arxiv.org/abs/2512.08848v2. Full licence notice: \texttt{licenses/arxiv-2512.08848-CC-BY-4.0.txt}.\par\smallskip

\noindent\textit{Editorial scope.} Counted unit: the article's lemma lemma:canonicalactionrightboundedoperators (source0.tex lines 1168--1171). The article's complete original proof of it is delivered with it (source0.tex lines 1172--1212, one \texttt{proof} environment of 41 lines). No mathematics is written, repaired or re-derived: the statement and the proof are the article's own lines, in the article's own notation, and no retained line is substituted or re-flowed.  No further source-local context is needed: every cross-reference of every retained line resolves inside the two delivered spans.  Nothing was omitted from any retained span, so the package carries no omission record.  No retained line contains a citation, so the wrapper carries no bibliography and no bibliography entry.  The retained lines contain the article's own diagram code, which the wrapper typesets with the standard tikz packages; no image file is delivered and the package declares no asset dependency.  Editorial formatting only, in addition: an AMS \texttt{amsart} preamble that restates the article's own declarations and macros used by the retained lines, namely the environments \texttt{lemma} on separate counters, together with the standard packages those lines need.  The retained environments print in this document as Lemma 1 in order of appearance, whereas the article prints the same items with the numbering of its own full document.  The complete native source of the article as distributed in the arXiv e-print for this exact version is bundled unchanged as \texttt{sources/190697/source0.tex}.
\begin{lemma}
    Given $V \in \Vect(\cC)$, there is a canonical left action of $\cL_V^l$ on $V$. If $H \in \Hilb(\cC)$, there is a canonical left action of $\bB^l_H$ on $\For(H)$.
    \label{lemma:canonicalactionrightboundedoperators}
\end{lemma}

\begin{proof}
The proof of the first claim is identitical to the proof of the second. We only explain the latter. For simplicity, we shall ommit the notation for the forgetful functor $\For: \Hilb(\cC) \to \Vect(\cC)$. There is a canonical linear map
\[ \bB^l_H(a_1) \odot H(a_2) \to H(a_1 \otimes a_2) . \]
Indeed, writing $\bB^l_H(a_1) \odot H(a_2) = \Hilb(\cC)(a_1 \otimes H,H) \odot \Hilb(\cC)(a_2,H)$ and $H(a_1 \otimes a_2) = \Hilb(\cC)(a_1 \otimes a_2, H)$, the above map is given by
    \[ f \odot g \mapsto f \circ (\id_{a_1} \otimes g) \ . \]
    Consider the composition
    \begin{align*}
    \cC(a, a_1 \otimes a_2) \odot \bB^l_H(a_1) \odot H(a_2) &\to \cC(a,a_1 \otimes a_2) \odot H(a_1 \otimes a_2)\to H(a) \\
    v \odot f \odot g &\mapsto v \odot f \circ (\id_{a_1} \otimes g)  \mapsto H(v) \left( f \circ (\id_{a_1} \otimes g) \right). 
    \end{align*}
    Now, since
    \[ (\bB^l_H \odot H) (a) \simeq \bigoplus_{a_1,a_2} \cC(a,a_1 \otimes a_2) \odot \bB^l_H(a_1) \odot H(a_2) \ , \]
    the maps defined above assemble into a morphism $\bB^l_H \odot H \to H$ in $\Vect(\cC)$.

    Using the generic notation $\triangleright: \bB^l_H \odot H \to H$ for the morphism defined above, and denoting by $\mu: \bB^l_H \odot \bB^l_H \to \bB^l_H$ the multiplication morphism for simplicity, we want to show that the diagram
    \begin{center}
        % https://tikzcd.yichuanshen.de/#N4Igdg9gJgpgziAXAbVABwnAlgFyxMJZABgBpiBdUkANwEMAbAVxiRAB12AjAIQD0ATgH0AEgAJO0CDgnd+w8ZKjSxIkAF9S6TLnyEUAJnJVajFm069Bo2VJlrN27HgJEjBk-WatEIB1pAMZz0iMg9qL3NfS3kbJRUHExgoAHN4IlAAMwEIAFskMhAcCCQARgizHw52XKZbZRlOLChRKA0A7Lyy6mKkI1NvC3YcASw6MBSGGFGUgAscdqyc-MQAZh6SxH6GOi4YBgAFHRd9EBn5kArB6OHR8cnprDmFxxBOlcLeteodvcPjkK+c4LK5RarNITAGLWETqeoqJotRLqIA
\begin{tikzcd}
\bB^l_H \odot \bB^l_H \odot H \arrow[rr, "\mu \odot \id_H"] \arrow[dd, "\id_{\bB^l_H} \odot \triangleright"'] &  & \bB^l_H \odot H \arrow[dd, "\triangleright"] \\
                                                                                                      &  &                                              \\
\bB^l_H \odot H \arrow[rr, "\triangleright"']                                                         &  & H                                           
\end{tikzcd}
    \end{center}
    commutes. Analyzing this diagram component-wisely, one concludes that its commmutativity is equivalent to the commutativity of the diagram
    \begin{center}
    % https://tikzcd.yichuanshen.de/#N4Igdg9gJgpgziAXAbVABwnAlgFyxMJZABgBpiBdUkANwEMAbAVxiRAB12AJLBgIwD6wTgGMAwgF8AFHQEBGAAScIeALbwFXUlwCUS9tBX6e-IaMkyBAJn0qs6uJtKa9yqEc4nBw9uOmyAZm0dEAlSdExcfEIUK3IqWkYWNk9eb3N-eVkbZTUNLV1bdxxjNLNfC0Dg0PCQDGw8AiI4qwT6ZlZEDm4ynz9LOWyq3RqIhuiiMlbqduSu1NM+yvlgoo8exYzLK2GQsLGoppQAFniZpM6QADM1koBzW4UAC1G6yMaY5AA2M8SOthuoiwACcRAopJwsFAhNkJLY8o47q4DMVnq96odPj9pn85td9CIQWCIewoTDrHDcvYNEjwZDocBZHJKQYEfoyYyKXtahiPkRTjjZpdASijFIHkDQXTSQzYfDqY4nnsEjAoHd4ERQFdgRBVEgyCAcBAkHJzv95uxVExHvSBFwoK9tbqTdQjUg4rjLpwcMCsHQwHcGDBfXcnjhHTq9YggobjYgPQw6HwYAwAArvCZdENhkBmvHe33+wPBrCh8P7a6R-WuuMxxPJtMZo4gbPhvNemXlPgAIQAesC7SzDCVbVwI87EKdY0gAKyuui8NiqOhoOBuitOqNz6eIH6GhcMJcrtfGjdVyc1pAAdnni66y9X69qm+vl93t8P9+P64oEiAA
\begin{tikzcd}
{\Hilb(\cC)(a_1 \otimes H,H) \odot \Hilb(\cC)(a_2 \otimes H, H) \odot \Hilb(\cC)(a_3,H)} \arrow[rr, "\mu \odot \id_Hd"] \arrow[dd, "\id_{\bB^l_H} \odot \triangleright"'] &  & {\Hilb(\cC)(a_1a_2 \otimes H,H) \odot \Hilb(\cC)(a_3,H)} \arrow[dd, "\triangleright"] &  & f \odot g \odot h \arrow[rr, maps to] \arrow[dd, maps to]  &  & f \circ (\id_{a_2} \otimes g) \odot h \arrow[dd, maps to]   \\
                                                                                                                                                                    &  &                                                                                         &  &                                                            &  &                                                             \\
{\Hilb(\cC)(a_1,H) \odot \Hilb(\cC)(a_2a_3,H)} \arrow[rr, "\triangleright"']                                                                                      &  & {\Hilb(\cC)(a_1a_2a_3,H)}                                                              &  & f \odot (g \circ (\id_{a_2} \otimes h) \arrow[rr, maps to] &  & f \circ (\id_{a_2} \otimes g) (\id_{a_1} \otimes \id_{a_2})
\end{tikzcd}  
    \end{center}
which can be checked directly:
    \begin{center}
        \begin{tikzcd}
            f \odot g \odot h \arrow[rr, maps to] \arrow[dd, maps to]  &  & f \circ (\id_{a_1} \otimes g) \odot h \arrow[dd, maps to] \\
            & & \\
            f \odot (g \circ (\id_{a_2} \otimes h) \arrow[rr, maps to] &  & f \circ (\id_{a_2} \otimes g) \circ (\id_{a_1} \otimes \id_{a_2} \otimes h)
        \end{tikzcd}
    \end{center}
\end{proof}

\end{document}
